\documentclass[10pt,a4paper]{ctexart}
\usepackage[margin=2cm]{geometry}
\usepackage{graphicx}
\usepackage{float}      % 支持 [H] 浮动体选项
\usepackage{hyperref}
\usepackage{amsmath}    % 提供方程环境
\usepackage{physics}    % 简化向量、微分算子语法
\newcommand{\myspace}[1]{\par\vspace{#1\baselineskip}} %自定义空一行
\usepackage{fancyhdr}
\pagestyle{fancy}
\fancyhf{}  % 清除所有页眉页脚
\fancyfoot[C]{\thepage}  % 在页脚中央设置页码
\renewcommand{\headrulewidth}{0pt}  % 去掉页眉横线
\renewcommand{\footrulewidth}{0pt}  % 去掉页脚横线

\usepackage{hyperref}
\hypersetup{
    colorlinks=true,
    linkcolor=blue,
    filecolor=magenta,      
    urlcolor=cyan,
    pdftitle={Overleaf Example},
    pdfpagemode=FullScreen,
    }

\usepackage{titlesec}

% 重新定义section格式
\titleformat{\section}
  {\normalfont\Large\bfseries}
  {\chinese{section}、}
  {0pt}
  {}

% 给出常用字号的自定义指令
\newcommand{\chuhao}{\fontsize{42pt}{48pt}\selectfont}
\newcommand{\xiaochu}{\fontsize{36pt}{42pt}\selectfont}
\newcommand{\xiaoer}{\fontsize{18pt}{21.6pt}\selectfont}
\newcommand{\sanhao}{\fontsize{16pt}{20pt}\selectfont}
\newcommand{\xiaosan}{\fontsize{15pt}{18pt}\selectfont}
\newcommand{\sihao}{\fontsize{14pt}{18pt}\selectfont}
\newcommand{\xiaosi}{\fontsize{12pt}{16pt}\selectfont}
\newcommand{\wuhao}{\fontsize{10.5pt}{12.6pt}\selectfont}

% 定义中文数字转换
\titleformat{\section}
  {\normalfont\Large\bfseries}
  {\zhnum{section}、}  % 使用\zhnum而不是\chinese
  {0pt}
  {}

% 标题设置
\title{\xiaochu\textbf{物理实验预习报告}}
\author{}
\date{}

\begin{document}

\vspace*{36pt} % 顶部空白

% 居中填写区域
\begin{center}   
    \xiaochu\textbf{物理实验预习报告}
    % 预留空白区域
    \vspace{13pt}
    \vspace{13pt}
    \vspace{13pt}
    \vspace{13pt}
    
    \sanhao\textbf{实验名称：\underline{\makebox[10cm]{ 密里根油滴实验}}} \\[10pt]
    \sanhao\textbf{指导教师：\underline{\makebox[10cm]{ 朱韵嫱 }}} \\[30pt]
    
    % 预留空白区域
    \myspace{6}
    
   \sihao\textbf{班级：\underline{\makebox[6cm]{}}} \\[10pt]
    \sihao\textbf{姓名：\underline{\makebox[6cm]{}}} \\[10pt]
    \sihao\textbf{学号：\underline{\makebox[6cm]{}}} \\[30pt]
    
   \sihao\textbf{实验日期: 
   \underline{\makebox[0.8cm]{2025}}  年 
   \underline{\makebox[0.4cm]{10}}月
   \underline{\makebox[0.4cm]{29}}日  
   星期\underline{\makebox[0.6cm]{三}}上午}
    
    \vspace{18pt}
    \sihao 浙江大学物理实验教学中心
    
\end{center}

\newpage

    \section{实验综述}
    
    \xiaosi （自述实验现象、实验原理和实验方法，不超过300字，5分）

本实验通过观察油滴在电场中的运动来测量基本电荷量。首先使用手机拍摄油滴的布朗运动现象，然后选择合适的油滴进行测量。
实验采用匀速下降法测定油滴的质量m， 具体方程如下

$m=\frac{4}{3}\pi r^3 \rho, 6\pi r \eta v = mg , v = \frac{l}{t}$。

接着通过静态平衡法测量油滴的带电量q,

$mg=q\frac{V}{d}$。

最终通过数据处理确定基本电荷e的数值。
实验过程中需要观察油滴的运动状态，并记录相关数据用于后续计算。
    
\myspace{2}
    
    \section{实验重点}
    
    \xiaosi（简述本实验的学习重点，不超过100字，3分）

实验的关键在于选择合适的油滴进行测量，油滴大小要适中，运动速度要适宜。
需要熟练掌握匀速下降法和静态平衡法的操作步骤，准确测量油滴的运动时间和平衡电压。
数据处理时要正确运用公式计算油滴的质量和带电量，最终通过分析多组数据确定基本电荷的数值。
    \myspace{2}
    
    \section{实验难点}

    \xiaosi（简述本实验的实现难点，不超过100字，2分）

实验的主要难点在于油滴的选择和观测，需要耐心寻找大小合适的油滴。
电场强度的精确调节和油滴运动时间的准确测量也是关键。
在数据处理阶段，需要正确处理测量误差，通过多组数据的统计分析来减小误差，最终得到准确的基本电荷值。
\myspace{2}


\newpage

\sihao\textbf{注意事项：}
\xiaosi
\begin{enumerate}
    \item 用PDF格式上传“预习报告”，文件名：学生姓名+学号+实验名称+周次。
    \item “预习报告”必须递交在“学在浙大”的本课程的对应实验项目的“作业”模块内。
    \item “预习报告”还须拷贝到“实验报告”中（便以教师批改）。
    \item 大学物理实验”课程选做实验使用本“预习报告”；必做实验无须写预习报告，前往“学在浙大”完成预习测试即可。
\end{enumerate}

\vspace{1cm}
\xiaosi\centering \textbf{浙江大学物理实验教学中心制}

\end{document}
